\documentclass[11pt]{article}
\usepackage{amsmath,amssymb,amsthm}

\newtheorem{theorem}{Theorem}
\newtheorem{lemma}[theorem]{Lemma}
\newtheorem{corollary}[theorem]{Corollary}
\newtheorem{remark}[theorem]{Remark}

\title{Explicit block-matrix unitary extensions of commuting isometries}
\author{}
\date{}

\begin{document}
\maketitle

\noindent
Let \(T,T'\in B(H)\) be isometries with \(TT'=T'T\).
We give \emph{explicit block operator matrices} for commuting unitary extensions \(U,U'\).

\bigskip
%% ------------------------------------------------------------------
\section*{0. Notation}
Let
\[
\mathcal{D}_*=\ker T^*=H\ominus TH,\qquad
\mathcal{D}'_*=\ker{T'}^*=H\ominus T'H.
\]
Write \(\ell^2(\mathbb{N},\mathcal{E})=\bigoplus_{n=1}^\infty\mathcal{E}\) and
\(\ell^2(\mathbb{Z},\mathcal{E})=\bigoplus_{n\in\mathbb{Z}}\mathcal{E}\).
Let \(S_{\mathcal{E}}\) denote the unilateral shift on \(\ell^2(\mathbb{N},\mathcal{E})\), so
\[
S_{\mathcal{E}}(e_1,e_2,e_3,\dots)=(0,e_1,e_2,\dots),\qquad
S_{\mathcal{E}}^*(e_1,e_2,e_3,\dots)=(e_2,e_3,\dots).
\]
Let \(J_{\mathcal{E}}:\mathcal{E}\hookrightarrow H\) be the inclusion (when \(\mathcal{E}\subset H\)), and define
\[
E_{\mathcal{E}}:\ell^2(\mathbb{N},\mathcal{E})\to H,\qquad
E_{\mathcal{E}}(e_1,e_2,e_3,\dots)=J_{\mathcal{E}}e_1.
\]
Thus \(E_{\mathcal{E}}^*h=(P_{\mathcal{E}}h,0,0,\dots)\).

\bigskip
%% ------------------------------------------------------------------
\section*{1. Unitary extension of a single isometry (block \(2\times 2\))}

\begin{lemma}
Let \(V\in B(H)\) be an isometry and set \(\mathcal{D}=\ker V^*\).
On the space
\[
K_V=\ell^2(\mathbb{N},\mathcal{D})\oplus H,
\]
the block operator matrix
\begin{equation}
\label{eq:single-U}
U_V
=
\begin{pmatrix}
S_{\mathcal{D}}^* & 0 \\
E_{\mathcal{D}} & V
\end{pmatrix}
\end{equation}
is a unitary extension of \(V\): \(U_V\big|_H=V\). Minimality means
\(K_V=\bigvee_{n\in\mathbb{Z}}U_V^n H\).
\end{lemma}

\begin{proof}
For \((x,h)\in \ell^2(\mathbb{N},\mathcal{D})\oplus H\),
\[
U_V\begin{pmatrix}x\\ h\end{pmatrix}
=
\begin{pmatrix}S_{\mathcal{D}}^*x\\ Vh+E_{\mathcal{D}}x\end{pmatrix}.
\]
Since \(E_{\mathcal{D}}x\in\mathcal{D}=H\ominus VH\),
\begin{align*}
\bigl\|U_V(x,h)\bigr\|^2
&=\|S_{\mathcal{D}}^*x\|^2+\|Vh\|^2+\|E_{\mathcal{D}}x\|^2
=\bigl(\|x\|^2-\|e_1\|^2\bigr)+\|h\|^2+\|e_1\|^2
=\|x\|^2+\|h\|^2.
\end{align*}
So \(U_V\) is isometric. Given \((y,k)\), set \(x=(P_{\mathcal{D}}k,y_1,y_2,\dots)\) and
\(h=V^*k\). Then \(S_{\mathcal{D}}^*x=y\) and \(Vh+E_{\mathcal{D}}x=k\), so \(U_V\) is
surjective, hence unitary. Clearly \(U_V(0,h)=(0,Vh)\).
\end{proof}

\noindent
In particular, the unitary extension of \(T\) is the block matrix
\begin{equation}
\label{eq:U-of-T}
U
=
\begin{pmatrix}
S_{\mathcal{D}_*}^* & 0 \\
E_{\mathcal{D}_*} & T
\end{pmatrix}
\quad\text{on}\quad
K_1=\ell^2(\mathbb{N},\mathcal{D}_*)\oplus H.
\end{equation}

\bigskip
%% ------------------------------------------------------------------
\section*{2. The commuting extension \(\widetilde{T}'\) on \(K_1\) (block form)}

Identify the standard orthonormal ``past'' basis vectors with negative powers of \(U\):
\[
U^{-n}d=(0,\dots,0,d,0,\dots;0)\in K_1
\qquad(d\in\mathcal{D}_*,\ n\ge 1).
\]
The inductive formula \(\widetilde{T}'(U^n h)=U^n T'h\) (\(n\in\mathbb{Z}\), \(h\in H\)) yields an
isometry \(\widetilde{T}'\) on \(K_1\) with \(\widetilde{T}'\big|_H=T'\) and
\(\widetilde{T}'U=U\widetilde{T}'\).

When \(T,T'\) \emph{doubly} commute (\(T^*T'=T'T^*\)), one has
\(T'\mathcal{D}_*\subset\mathcal{D}_*\), so \(\Gamma:=T'|_{\mathcal{D}_*}\) is an isometry and
\begin{equation}
\label{eq:double-block}
\widetilde{T}'
=
\begin{pmatrix}
I_{\ell^2(\mathbb{N})}\otimes\Gamma & 0 \\
0 & T'
\end{pmatrix}
\quad\text{on}\quad
\ell^2(\mathbb{N},\mathcal{D}_*)\oplus H.
\end{equation}
In the general (not doubly commuting) case, \(\widetilde{T}'\) mixes the past with \(H\), and the
cleanest fully explicit matrices for \emph{both} unitaries simultaneously are the Berger--Coburn--Lebow
Laurent matrices of Section~3.

\medskip
If \(\widetilde{T}'\) is not yet unitary, apply~\eqref{eq:single-U} to \(\widetilde{T}'\) on \(K_1\):
with \(\mathcal{D}_{\sim}=\ker\widetilde{T}'{}^*\),
\begin{equation}
\label{eq:Uprime}
U'
=
\begin{pmatrix}
S_{\mathcal{D}_{\sim}}^* & 0 \\
E_{\mathcal{D}_{\sim}} & \widetilde{T}'
\end{pmatrix}
\quad\text{on}\quad
K=\ell^2(\mathbb{N},\mathcal{D}_{\sim})\oplus K_1,
\end{equation}
and extend \(U\) to a unitary \(\widetilde{U}\) on \(K\) by the same inductive rule
(replace \((T,T')\) by \((\widetilde{T}',U)\)). Then \(\widetilde{U},U'\) are commuting unitaries
extending \(T,T'\).

\bigskip
%% ------------------------------------------------------------------
\section*{3. Fully explicit Laurent matrices (pure case): BCL model}

Assume the product \(TT'\) is a \emph{pure} isometry (equivalently
\((TT')^{*n}\to 0\) strongly). Set
\[
\mathcal{W}=\ker(TT')^*=H\ominus TT'H.
\]
Let
\[
W_1=\ker T^*=\mathcal{D}_*,\qquad
W_2=\ker{T'}^*=\mathcal{D}'.
\]
One has the orthogonal decompositions
\[
\mathcal{W}=W_1\oplus T'W_1=W_2\oplus TW_2
\]
(up to the usual identification of wandering subspaces for a pure pair). Define a unitary
\(U_{\mathcal{W}}\) on \(\mathcal{W}\) and an orthogonal projection \(P\) on \(\mathcal{W}\) by the
Berger--Coburn--Lebow formulae
\begin{equation}
\label{eq:UW-P}
U_{\mathcal{W}}
=
\begin{pmatrix}
T'\big|_{W_1} & 0 \\
0 & T^*\big|_{TW_2}
\end{pmatrix}
:
W_1\oplus TW_2
\;\longrightarrow\;
T'W_1\oplus W_2,
\qquad
P=P_{W_2}.
\end{equation}
(Other equivalent normalizations replace \(P_{W_2}\) by \(P_{W_1}\).)

\begin{theorem}[BCL unitary extensions]
Identify \(H\) with the Hardy space \(H^2(\mathcal{W})\subset L^2(\mathbb{T};\mathcal{W})\), so that
\[
T\simeq M_{\Phi},\quad
T'\simeq M_{\Psi}
\quad\text{on }H^2(\mathcal{W}),
\]
where the \(\mathrm{B}(\mathcal{W})\)-valued inner functions
\begin{equation}
\label{eq:PhiPsi}
\Phi(z)=U_{\mathcal{W}}^*\bigl(P+zP^\perp\bigr),
\qquad
\Psi(z)=\bigl(P^\perp+zP\bigr)U_{\mathcal{W}}
\end{equation}
are \emph{unitary-valued for a.e.\ \(z\in\mathbb{T}\)}. Consequently the multiplication operators
\begin{equation}
\label{eq:UU-L2}
U=M_{\Phi},\qquad U'=M_{\Psi}
\quad\text{on}\quad
K=L^2(\mathbb{T};\mathcal{W})
\end{equation}
are commuting unitaries, and
\[
U\big|_{H^2(\mathcal{W})}=T,\qquad
U'\big|_{H^2(\mathcal{W})}=T'.
\]
\end{theorem}

\subsection*{Laurent block matrices on \(\ell^2(\mathbb{Z},\mathcal{W})\)}

Under the Fourier identification \(L^2(\mathbb{T};\mathcal{W})\cong\ell^2(\mathbb{Z},\mathcal{W})\),
write vectors \(\xi=(\xi_n)_{n\in\mathbb{Z}}\), \(\xi_n\in\mathcal{W}\). Then~\eqref{eq:PhiPsi}
become the bi-infinite Laurent matrices
\begin{equation}
\label{eq:Laurent-U}
\boxed{
\begin{aligned}
(U\xi)_n
&=
U_{\mathcal{W}}^*P\,\xi_n
+
U_{\mathcal{W}}^*P^\perp\,\xi_{n-1},
\\[1mm]
(U'\xi)_n
&=
P^\perp U_{\mathcal{W}}\,\xi_n
+
P\,U_{\mathcal{W}}\,\xi_{n-1}.
\end{aligned}
}
\end{equation}
Equivalently, as operator matrices with respect to \(\bigoplus_{n\in\mathbb{Z}}\mathcal{W}\),
\begin{equation}
\label{eq:matrix-form}
U
=
\begin{pmatrix}
\ddots & & & & \\
& U_{\mathcal{W}}^*P & 0 & & \\
& U_{\mathcal{W}}^*P^\perp & U_{\mathcal{W}}^*P & 0 & \\
& 0 & U_{\mathcal{W}}^*P^\perp & U_{\mathcal{W}}^*P & \\
& & & & \ddots
\end{pmatrix},
\qquad
U'
=
\begin{pmatrix}
\ddots & & & & \\
& P^\perp U_{\mathcal{W}} & 0 & & \\
& P\,U_{\mathcal{W}} & P^\perp U_{\mathcal{W}} & 0 & \\
& 0 & P\,U_{\mathcal{W}} & P^\perp U_{\mathcal{W}} & \\
& & & & \ddots
\end{pmatrix}.
\end{equation}
The subspace \(H\simeq H^2(\mathcal{W})\) is \(\bigoplus_{n\ge 0}\mathcal{W}\), which is invariant under
both \(U\) and \(U'\), and the compressions / restrictions recover \(T,T'\).

\begin{proof}[Proof that \(U,U'\) are commuting unitaries extending \(T,T'\)]
For a.e.\ \(z\in\mathbb{T}\), \(P+zP^\perp\) is unitary on \(\mathcal{W}\) (it is a
symmetry times a unimodular scalar on the two reducing subspaces \(\operatorname{Ran}P\) and
\(\operatorname{Ran}P^\perp\)). Since \(U_{\mathcal{W}}\) is unitary, both \(\Phi(z)\) and
\(\Psi(z)\) are unitary. Hence \(M_{\Phi}\) and \(M_{\Psi}\) are unitary on
\(L^2(\mathbb{T};\mathcal{W})\). They commute because the \(\mathrm{B}(\mathcal{W})\)-valued
functions \(\Phi,\Psi\) commute pointwise:
\[
\Phi(z)\Psi(z)=U_{\mathcal{W}}^*(P+zP^\perp)(P^\perp+zP)U_{\mathcal{W}}
=z\,I
=\Psi(z)\Phi(z).
\]
Restriction to \(H^2(\mathcal{W})\) yields commuting isometries unitarily equivalent to
\((T,T')\) by the BCL theorem. The formulae~\eqref{eq:Laurent-U}--\eqref{eq:matrix-form}
are the Fourier expansions of multiplication by~\eqref{eq:PhiPsi}.
\end{proof}

\bigskip
%% ------------------------------------------------------------------
\section*{4. Non-pure case}

In general, let \(H_u=\bigcap_{n\ge 0}(TT')^n H\) be the maximal subspace reducing \(TT'\) to a
unitary (Suciu / BCL). Then \(H=H_u\oplus H_s\) reduces both \(T\) and \(T'\), with
\((T|_{H_u},T'|_{H_u})\) already unitary, and \((T|_{H_s},T'|_{H_s})\) pure. Apply
Section~3 on \(H_s\) and take the direct sum with the unitary pair on \(H_u\):
\[
U=T|_{H_u}\oplus M_{\Phi},
\qquad
U'=T'|_{H_u}\oplus M_{\Psi}
\quad\text{on}\quad
K=H_u\oplus L^2(\mathbb{T};\mathcal{W}).
\]

\bigskip
%% ------------------------------------------------------------------
\section*{5. Summary: the matrices to copy}

\begin{itemize}
\item \textbf{One isometry.} On \(\ell^2(\mathbb{N},\ker V^*)\oplus H\),
\[
U_V=\begin{pmatrix}S^*&0\\ E&V\end{pmatrix}.
\]

\item \textbf{Commuting pair, pure product (main explicit form).} On
\(\ell^2(\mathbb{Z},\mathcal{W})\), with \(U_{\mathcal{W}},P\) from~\eqref{eq:UW-P},
\[
(U\xi)_n=U_{\mathcal{W}}^*P\,\xi_n+U_{\mathcal{W}}^*P^\perp\,\xi_{n-1},
\qquad
(U'\xi)_n=P^\perp U_{\mathcal{W}}\,\xi_n+P\,U_{\mathcal{W}}\,\xi_{n-1}.
\]

\item \textbf{Doubly commuting special case.} On
\(\ell^2(\mathbb{N},\mathcal{D}_*)\oplus H\),
\[
U=\begin{pmatrix}S_{\mathcal{D}_*}^*&0\\ E_{\mathcal{D}_*}&T\end{pmatrix},
\qquad
\widetilde{T}'=\begin{pmatrix}I\otimes T'|_{\mathcal{D}_*}&0\\ 0&T'\end{pmatrix},
\]
then unitarize \(\widetilde{T}'\) by~\eqref{eq:single-U} and extend \(U\) in parallel.
\end{itemize}

\end{document}
